% Generated draft: requires mathematical and source review.
\documentclass[11pt]{article}
\usepackage[margin=25mm]{geometry}
\usepackage{fontspec}
\setmainfont{Latin Modern Roman}
\usepackage{amsmath,amssymb,mathtools,unicode-math}
\setmathfont{Latin Modern Math}
\usepackage{enumitem,needspace,xurl,hyperref}
\setlength{\parindent}{0pt}
\setlength{\parskip}{5pt}
\setlength{\emergencystretch}{3em}
\newcommand{\N}{\mathbb{N}}
\newcommand{\Z}{\mathbb{Z}}
\newcommand{\Q}{\mathbb{Q}}
\newcommand{\R}{\mathbb{R}}
\newcommand{\C}{\mathbb{C}}
\newcommand{\F}{\mathbb{F}}
\newcommand{\A}{\mathbb{A}}
\newcommand{\PP}{\mathbb{P}}
\newcommand{\E}{\mathbb{E}}
\newcommand{\eps}{\varepsilon}
\newcommand{\dd}{\,\mathrm{d}}
\begin{document}
\section*{UnsolvedMath ID: 701159}
\textbf{Source problem code:} catalogue entry
\subsection{Definitions and mathematical statement}
For a finite group \(G\), let \(\operatorname{Irr}(G)\) denote its ordinary irreducible complex characters, and let \(\chi(1)\in\Z_{\geq 1}\) be the degree of \(\chi\). This uses the standard ordinary-character interpretation; the supplied statement does not explicitly specify a character field. For a positive integer \(n\), define
\[
\pi(n)=\{p\in\Z_{\geq 2}:p\text{ is prime and }p\mid n\},
\qquad
\rho(G)=\bigcup_{\chi\in\operatorname{Irr}(G)}\pi(\chi(1)).
\]
In particular, \(\pi(1)=\varnothing\). The question is whether every finite solvable group \(G\) has a character \(\chi\in\operatorname{Irr}(G)\) such that
\[
\bigl|\pi(\chi(1))\bigr|\geq\frac{|\rho(G)|}{2},
\]
or, equivalently, such that \(\bigl|\pi(\chi(1))\bigr|\geq\lceil |\rho(G)|/2\rceil\). Distinct prime divisors are counted without multiplicity, and no exclusion is made when \(\rho(G)=\varnothing\).
\subsection{Short English statement}
Must every finite solvable group have an ordinary irreducible complex character whose degree is divisible by at least half of all the distinct primes occurring in its irreducible character degrees? The supplied registry labels this question open; that label is not independently verified here.
\subsection{Sources}
\begin{itemize}
\item Supplied local UnsolvedMath registry record: permanent ID \(701159\), exact source problem code \texttt{catalogue entry}, title \emph{Kourovka Notebook Problem 15.71}; upstream revision \texttt{d371ee917ac9647fbb06cd6fc529deec6f3dda82}.
\item Supplied bibliographic source: \emph{The Kourovka Notebook (Unsolved Problems in Group Theory)}, 15th issue (2002), Problem 15.71; also identified in the record with the 21st issue edition (2026), \texttt{arXiv:1401.0300v48}. These bibliographic details are supplied attribution and have not been independently checked.
\item The statement credits B. Huppert parenthetically. The registry identifies P. P. P\'alfy as proposer, with proposed year 2002. Both attributions are retained without inferring their precise relationship.
\item The supplied collection description names E. I. Khukhro and V. D. Mazurov as editors of the Kourovka Notebook collection covering Issues 1--21 (1965--2026).
\end{itemize}
\end{document}
